
\chapter{Contexte Général du Projet}
\section*{Sommaire}
\begin{quote}
\small
\textbf{1} \hspace{1em} Introduction \dotfill \pageref{sec:intro_ch1}\\
\textbf{2} \hspace{1em} Présentation de l'organisme d'accueil \dotfill \pageref{sec:entreprise}\\
\textbf{3} \hspace{1em} Présentation du projet \dotfill \pageref{sec:projet}\\
\hspace{2.5em} 3.1 \hspace{0.5em} Contexte du projet \dotfill \pageref{subsec:contexte}\\
\hspace{2.5em} 3.2 \hspace{0.5em} Problématique \dotfill \pageref{subsec:problematique}\\
\hspace{2.5em} 3.3 \hspace{0.5em} Objectifs du projet \dotfill \pageref{subsec:objectifs}\\
\hspace{2.5em} 3.4 \hspace{0.5em} Périmètre du projet \dotfill \pageref{subsec:perimetre}\\
\hspace{2.5em} 3.5 \hspace{0.5em} Outil \dotfill \pageref{subsec:qqoqcp}\\
\textbf{4} \hspace{1em} Organisation et planification du stage \dotfill \pageref{sec:organisation}\\
\hspace{2.5em} 4.1 \hspace{0.5em} Méthodologie de travail \dotfill \pageref{subsec:methodo}\\
\hspace{2.5em} 4.2 \hspace{0.5em} Planning général \dotfill \pageref{subsec:planning}\\
\hspace{2.5em} 4.3 \hspace{0.5em} Livrables attendus \dotfill \pageref{subsec:livrables}\\
\textbf{5} \hspace{1em} Conclusion \dotfill \pageref{sec:conclu_ch1}\\
\end{quote}

\newpage

% ------------------------------------------------------------
\section{Introduction}
\label{sec:intro_ch1}

Dans un contexte industriel où les systèmes embarqués communicants
occupent une place de plus en plus centrale dans la gestion des flottes
de véhicules et la collecte de données terrain, la conception d'un
firmware robuste, générique et maintenable représente un enjeu
stratégique majeur pour les entreprises du secteur télématique. C'est
précisément dans ce contexte que s'inscrit le présent projet de fin
d'études, réalisé au sein de Smart Automation Technologies ,
entreprise marocaine spécialisée dans les solutions embarquées et
l'intelligence artificielle appliquée. Ce premier chapitre a pour
objectif de poser les bases conceptuelles et organisationnelles
indispensables à la compréhension globale du travail accompli. Il
présente successivement l'entreprise d'accueil, la problématique
technique qui a motivé ce stage, les objectifs fixés, le périmètre
du projet, ainsi que la méthodologie et la planification adoptées
tout au long des vingt semaines de développement.

% ------------------------------------------------------------
\section{Présentation de l'organisme d'accueil}
\label{sec:entreprise}

\begin{figure}[h]
    \centering
    
\includegraphics[width=0.35\textwidth]{figures/logo_SAT.jpeg}
    \caption{Logo de Smart Automation Technologies}
    \label{fig:logo_sat}
\end{figure}

Smart Automation Technologies (SAT) est une entreprise marocaine de
technologie basée à Tanger, fondée par un groupe de consultants seniors
aux profils pluridisciplinaires couvrant les systèmes d'information,
la science des données et l'intelligence artificielle embarquée.
Depuis sa création, SAT s'est positionnée comme un partenaire
technologique de référence pour les entreprises et administrations
cherchant à moderniser leurs systèmes d'information et à intégrer des
solutions intelligentes dans leurs processus opérationnels. L'entreprise
se distingue par une approche orientée résultats, qui privilégie la
compréhension approfondie des besoins métier de ses clients avant toute
proposition de solution technique. Cette démarche lui permet de livrer
des systèmes réellement adaptés aux contraintes opérationnelles,
budgétaires et humaines de chaque organisation. SAT intervient aussi
bien dans la conception initiale de systèmes complexes que dans leur
évolution, leur maintenance et leur optimisation sur le long terme,
assurant ainsi une relation de partenariat durable avec ses clients.

\paragraph{Domaines d'activité.}
Les activités de SAT s'articulent autour de quatre pôles technologiques
complémentaires. Le premier pôle concerne le développement de systèmes
embarqués et IoT, où l'entreprise conçoit des firmwares sur mesure, des
cartes électroniques dédiées et des solutions de connectivité pour des
applications industrielles, automobiles et de mobilité urbaine. Le
deuxième pôle porte sur l'intelligence artificielle et la science des
données, avec le développement de modèles prédictifs, d'algorithmes de
traitement en temps réel et de solutions d'analyse à grande échelle
déployées sur des plateformes cloud et edge. Le troisième pôle recouvre
le conseil et la gouvernance IT, incluant l'assistance à maîtrise
d'ouvrage, l'audit d'architecture logicielle et le benchmarking de
systèmes d'information. Enfin, le quatrième pôle est dédié au support
et à la formation, avec des programmes ciblés sur les outils et
technologies déployés chez les clients, garantissant une montée en
compétences progressive et durable des équipes.

\paragraph{Positionnement et valeurs.}
La vision stratégique de SAT repose sur trois piliers fondamentaux :
l'innovation technologique continue, la relation de proximité avec ses
clients et la valorisation du capital humain. L'entreprise cultive une
culture interne axée sur l'excellence technique, l'intégrité dans les
engagements pris, l'esprit d'équipe et l'engagement envers la qualité
des livrables. Ces valeurs, illustrées dans la
figure~\ref{fig:valeurs}, se traduisent concrètement dans les pratiques
quotidiennes : revues de code régulières, encadrement des stagiaires,
veille technologique active et participation à des projets de recherche
appliquée en partenariat avec des universités marocaines.

\begin{figure}[h]
    \centering
    \includegraphics[width=0.55\textwidth]{figures/valeurs_sat.jpeg}
    \caption{Valeurs fondamentales de Smart Automation Technologies}
    \label{fig:valeurs}
\end{figure}

\paragraph{Intérêt du projet firmware pour l'entreprise.}
Dans le domaine des systèmes embarqués et de la télématique, SAT
développe une gamme d'appareils déployés dans des flottes de véhicules
pour la collecte, le traitement et la transmission de données de
localisation, de diagnostic moteur et d'état des véhicules. Face à la
diversité croissante de ses références matérielles et à la complexité
grandissante de leur maintenance logicielle respective, SAT a identifié
la nécessité stratégique de disposer d'un firmware commun, générique
et configurable capable de couvrir l'ensemble de la gamme. Ce projet
de fin d'études répond directement à cet enjeu en proposant une
architecture unifiée, réduisant la dette technique accumulée et
accélérant les cycles de développement des futures variantes produits.

% ------------------------------------------------------------
\section{Présentation du projet}
\label{sec:projet}

\subsection{Contexte du projet}
\label{subsec:contexte}

La gamme d'appareils télématiques de SAT comprend plusieurs références
matérielles qui, bien que partageant un socle technologique commun basé
sur le microcontrôleur ESP32-P4, présentent des spécificités
fonctionnelles propres à chaque produit. Ces différences se manifestent
à plusieurs niveaux : interfaces périphériques distinctes selon les
variantes (GPS, GSM/LTE, accéléromètre, bus CAN, entrées/sorties
numériques), protocoles de communication variés selon les cas d'usage
(MQTT over TLS, HTTP, CAN/J1939), et exigences de gestion d'énergie
différentes selon que l'appareil est alimenté en permanence sur le
réseau de bord du véhicule ou sur batterie autonome. À ces contraintes
matérielles s'ajoutent des exigences systèmes sévères inhérentes aux
environnements embarqués communicants : ressources mémoire limitées
(quelques mégaoctets de RAM et de Flash), contraintes temps réel
strictes sur les tâches critiques (acquisition GPS, publication MQTT,
gestion de l'état d'énergie), et obligation de maintenance à distance
via des mécanismes de configuration runtime et de diagnostic sans
intervention physique sur les véhicules du parc déployé. Face à cette
diversité, l'équipe R\&D de SAT se retrouvait jusqu'ici contrainte de
maintenir autant de bases de code distinctes qu'il existe de variantes
produits, engendrant une dette technique croissante, des risques de
régression lors des mises à jour de sécurité et un allongement
significatif des cycles de mise sur le marché de nouveaux produits.

\subsection{Problématique}
\label{subsec:problematique}

La problématique centrale de ce projet émerge naturellement du constat
décrit ci-dessus. Maintenir des bases de code fragmentées pour chaque
variante matérielle est une approche qui atteint rapidement ses limites
à mesure que le nombre de produits croît et que les exigences de
sécurité et de fiabilité se durcissent. La question fondamentale à
laquelle ce projet cherche à répondre est la suivante :

\begin{quote}
\textit{Comment concevoir un firmware générique, modulaire et
configurable, capable de s'adapter à plusieurs configurations
matérielles et logicielles d'appareils télématiques, tout en
garantissant fiabilité, maintenabilité et évolutivité ?}
\end{quote}

Cette problématique soulève plusieurs défis techniques interdépendants
qui structurent l'ensemble de la démarche de conception. Le premier
défi concerne la généricité et la configurabilité : il s'agit
d'abstraire suffisamment les différences matérielles entre variantes
pour qu'un seul binaire puisse couvrir toute la gamme en s'appuyant
sur un registre de capacités 32 bits et une couche HAL bien définie.
Le deuxième défi porte sur la modularité : l'architecture doit
permettre l'ajout d'un nouveau module (capteur, protocole, service)
sans perturber les modules existants ni nécessiter de recompilation
globale. Le troisième défi concerne la concurrence et la gestion des
ressources sous FreeRTOS, où des dizaines de tâches aux priorités
hétérogènes doivent cohabiter dans un espace mémoire contraint tout
en respectant les budgets RAM et Stack de chaque composant. Enfin,
le quatrième défi touche à la sécurité et à la maintenance à distance,
avec l'intégration de mécanismes de mise à jour firmware Over-The-Air
sécurisés via TLS et de journalisation structurée permettant le
diagnostic à distance du parc déployé.

\subsection{Objectifs du projet}
\label{subsec:objectifs}

L'objectif principal de ce projet est de concevoir, implémenter et
valider une architecture firmware générique, modulaire et configurable
pour les appareils télématiques de SAT, basée sur l'ESP32-P4 sous
FreeRTOS et ESP-IDF v5.3.1, et livrable sous la forme d'une base de
code pérenne, documentée et immédiatement exploitable par l'équipe
R\&D. Pour atteindre cet objectif, plusieurs objectifs spécifiques
ont été définis et servent de jalons mesurables tout au long du
développement. Le premier objectif consiste à définir un modèle de
capacités matérielles et un modèle de configuration runtime permettant
à un seul binaire de s'adapter à n'importe quelle variante de la gamme
par simple lecture d'un registre persistant en NVS au démarrage. Le
deuxième objectif porte sur la mise en place d'une machine d'états
système hiérarchique à cinq états (\texttt{INIT}, \texttt{RUN},
\texttt{SLEEP}, \texttt{FAULT}, \texttt{UPDATE}) pilotant le cycle
de vie complet du firmware de manière déterministe et récupérable.
Les objectifs suivants couvrent l'implémentation des services core
--- configuration, stockage NVS, journalisation structurée, gestion
d'énergie et diagnostic interne --- ainsi que la création d'une
architecture modulaire avec HAL standardisé permettant l'ajout de
drivers par un développeur tiers sans risque de régression. Enfin,
les objectifs de validation comprennent la rédaction de tests
unitaires et d'intégration, le profilage RAM/Stack sur cible réelle
et la production d'une documentation technique de transfert.

\subsection{Périmètre du projet}
\label{subsec:perimetre}

\paragraph{Périmètre fonctionnel.}
Le projet couvre la conception et l'implémentation d'une architecture
firmware unifiée pour les trois variantes de la gamme SAT : Basic
(GPS uniquement), Pro (GPS, CAN, OTA, MQTT, Logger, Power avancé) et
Ultra Pro (toutes les fonctionnalités activées). Chaque variante est
définie par un registre de capacités 32 bits persisté en NVS, lu au
démarrage par le système d'initialisation pour configurer dynamiquement
les services et drivers à activer.

\paragraph{Périmètre firmware.}
Le développement porte sur les quatre couches de l'architecture :
HAL (abstraction matérielle), Drivers (implémentations concrètes des
périphériques), Services (logique applicative découplée) et Application
(orchestration globale et machine à états). L'ensemble est développé
sous ESP-IDF v5.3.1 avec FreeRTOS sur la plateforme ESP32-P4, en
utilisant Espressif-IDE sous Windows comme environnement de
développement principal.

\paragraph{Périmètre configuration et diagnostic.}
Les mécanismes de persistance NVS, de configuration runtime, de
journalisation structurée avec niveaux de sévérité et de diagnostic
interne sont inclus dans le périmètre. Le protocole de communication
UART avec le module de gestion OTA/MQTT développé par le collègue
spécialiste est également défini et implémenté côté firmware.

\paragraph{Limites du projet.}
Le développement du serveur OTA backend et de la plateforme MQTT
cloud est hors périmètre, ces composants étant pris en charge par
un collègue spécialiste en charge de l'infrastructure serveur. La
portabilité vers d'autres microcontrôleurs constitue une perspective
future et non un livrable de ce stage.

\subsection{Outil QQOQCP}
\label{subsec:qqoqcp}

L'outil QQOQCP (\textit{Qui~? Quoi~? Où~? Quand~? Comment~?
Pourquoi~?}) permet de synthétiser de manière structurée les aspects
essentiels du projet. Le tableau~\ref{tab:qqoqcp} ci-dessous présente
cette analyse synthétique, facilitant la compréhension du contexte
et des enjeux du projet Globalfleet.

\begin{table}[H]
    \centering
    \renewcommand{\arraystretch}{1.6}
    \caption{Outil QQOQCP du projet Globalfleet}
    \label{tab:qqoqcp}
    \begin{tabular}{|p{2cm}|p{10.5cm}|}
        \hline
        \rowcolor{gray!20}
        \textbf{Question} & \textbf{Réponse} \\
        \hline
        \textbf{Qui ?} &
        OULGHAZI Ihssane --- Stagiaire PFE Master Systèmes Embarqués
        et Robotique, FSTH, encadrée par M.~Jabir El Aaraj (SAT) et
        M.~Badr Benamrou (encadrant académique). \\
        \hline
        \textbf{Quoi ?} &
        Développement d'un firmware générique, modulaire et
        configurable pour microcontrôleur ESP32-P4 sous FreeRTOS,
        couvrant l'ensemble de la gamme d'appareils télématiques
        de SAT via une architecture en couches (HAL, Drivers,
        Services, Application), un Event Bus, un Ring Buffer et
        des services critiques (NVS, Power Management, MQTT/TLS). \\
        \hline
        \textbf{Où ?} &
        Smart Automation Technologies (SAT) --- Tanger, Maroc. \\
        \hline
        \textbf{Quand ?} &
        De Mars 2026 à Juillet 2026 (stage PFE de 6 mois). \\
        \hline
        \textbf{Comment ?} &
        \begin{itemize}[noitemsep, topsep=0pt, leftmargin=*]
            \item Audit et analyse critique des bases de code
                  firmware existantes.
            \item Conception de l'architecture modulaire en couches
                  sur ESP32-P4.
            \item Implémentation sous ESP-IDF v5.3.1 avec
                  Espressif-IDE, Git et CI/CD.
            \item Tests unitaires (pytest) et tests d'intégration
                  MQTT sur cible réelle.
            \item Profilage mémoire (heap monitor, stack watermark
                  FreeRTOS).
            \item Développement et validation sur les trois
                  variantes : Basic, Pro, Ultra Pro.
        \end{itemize} \\
        \hline
        \textbf{Pourquoi ?} &
        Résorber la dette technique de SAT en centralisant le code
        source de la gamme télématique dans un dépôt Git unifié,
        accélérer la mise sur le marché de nouveaux produits et
        disposer d'une plateforme firmware sécurisée, extensible
        et documentée à long terme. \\
        \hline
    \end{tabular}
\end{table}

% ------------------------------------------------------------
\section{Organisation et planification du stage}
\label{sec:organisation}

\subsection{Méthodologie de travail}
\label{subsec:methodo}

La méthodologie adoptée pour ce projet s'appuie sur une approche
progressive par phases, conçue pour s'adapter aux contraintes
spécifiques du développement firmware embarqué, où les décisions
d'architecture de bas niveau peuvent remettre en question des choix
effectués en phase amont. Plutôt que d'adopter une méthodologie Agile
stricte avec des sprints formalisés, l'approche retenue privilégie un
développement incrémental du noyau firmware avec validation
systématique de chaque module sur cible réelle avant toute intégration
dans l'architecture globale. Cette discipline de validation précoce
permet de détecter au plus tôt les problèmes de timing, de consommation
mémoire ou d'interférence entre tâches FreeRTOS, qui sont nettement
plus coûteux à corriger en phase d'intégration tardive. La gestion
de version est assurée via Git avec une chaîne CI de compilation
automatique permettant de vérifier en continu que l'ensemble des
22 composants ESP-IDF compile sans erreur après chaque modification.
Des points de revue réguliers avec l'encadrant professionnel permettent
de valider les livrables de chaque phase et d'ajuster les priorités
en fonction des contraintes techniques découvertes sur le terrain.

\subsection{Planning général}
\label{subsec:planning}

Le diagramme de Gantt présenté dans la figure~\ref{fig:gantt} illustre
la décomposition chronologique du projet en cinq phases principales,
s'étalant sur vingt semaines de stage. Le tableau~\ref{tab:planning}
synthétise ces phases avec leurs durées indicatives et leurs livrables
principaux.

\begin{table}[H]
    \centering
    \renewcommand{\arraystretch}{1.4}
    \caption{Planning synthétique du projet Globalfleet}
    \label{tab:planning}
    \begin{tabular}{|c|p{3.8cm}|p{3cm}|p{5cm}|}
        \hline
        \rowcolor{gray!20}
        \textbf{Phase} & \textbf{Intitulé} & \textbf{Durée} &
        \textbf{Livrables principaux} \\
        \hline
        1 & Étude du matériel et cadrage architecture
          & Semaines 1 – 3
          & Rapport d'analyse, choix technologiques documentés \\
        \hline
        2 & Squelette noyau RTOS
          & Semaines 3 – 6
          & Event Bus, Ring Buffer, gestion d'erreurs, CI compile \\
        \hline
        3 & Services core du firmware
          & Semaines 6 – 11
          & Config NVS, Storage, Logger, Power Manager, Diagnostic \\
        \hline
        4 & Modularité et intégration
          & Semaines 11 – 15
          & HAL complet, sélecteur de protocole, machine à états \\
        \hline
        5 & Tests, durcissement et documentation
          & Semaines 15 – 20
          & Rapport de tests, profilage RAM/Stack, documentation \\
        \hline
    \end{tabular}
\end{table}
Figure 1.3 présente le diagramme de Gantt actualisé du projet de stage, comparant l'avancement réel aux jalons initialement planifiés. Comme illustré dans la Fig. 1.3, les Phases 1 à 4 (S1 à S13), couvrant l'étude du matériel, le squelette du noyau RTOS, les services core du firmware ainsi que la modularité et l'intégration, sont entièrement achevées  de même que la tâche S14 relative aux tests unitaires et d'intégration. Seules les tâches S15 (profiling RAM/stack et release candidate) et S16 (préparation du transfert et documentation finale), encore en cours, ont été replanifiées sur la période du 3 au 7 juillet 2026. Cette représentation graphique met en évidence une avance d'environ un mois par rapport au calendrier prévisionnel initial, qui prévoyait l'achèvement de l'ensemble du projet à la fin du mois de juillet 2026.
\begin{figure}[H]
    \centering

\includegraphics[width=\textwidth]{figures/gantt.png}
    \caption{Diagramme de Gantt du projet Globalfleet}
    \label{fig:gantt}
\end{figure}

\subsection{Livrables attendus}
\label{subsec:livrables}

L'ensemble des livrables du projet Globalfleet peut être regroupé en
cinq catégories cohérentes avec les phases de développement. Les
livrables d'architecture comprennent le document d'architecture
firmware, la table des capacités et de configuration, ainsi que la
spécification de la machine d'états système à cinq états. Les
livrables d'infrastructure incluent le dépôt Git opérationnel avec
CI de compilation, l'Event Bus avec son interface module standardisée
et le catalogue exhaustif des codes d'erreurs du système. Les
livrables de services core couvrent les modules Config, Storage,
Logger, Power Manager et Diagnostic, entièrement implémentés, testés
et documentés avec leurs interfaces publiques. Les livrables
d'extensibilité comprennent le HAL complet avec son guide d'ajout
de driver et le sélecteur de protocole runtime permettant de basculer
dynamiquement entre les protocoles de communication disponibles.
Enfin, les livrables de validation et de documentation incluent le
rapport de tests unitaires et d'intégration, le rapport de profilage
RAM/Stack mesuré sur cible ESP32-P4 réelle, et la documentation
technique finale de transfert à destination de l'équipe R\&D de SAT.

% ------------------------------------------------------------
\section{Conclusion}
\label{sec:conclu_ch1}
Ce premier chapitre a permis de poser le cadre du projet : présentation de Smart Automation Technologies, exposé de la problématique — concevoir une architecture firmware générique, modulaire et configurable, adaptée à plusieurs variantes matérielles tout en garantissant performances temps réel, maintenabilité et sécurité — ainsi que la définition du périmètre, des objectifs et des livrables attendus. Le Chapitre 2 présentera l'état de l'art des technologies mobilisées (ESP32-P4, FreeRTOS, architectures modulaires HAL/Event-Driven/Services, protocoles de communication et sécurité) sur lesquelles s'appuie la solution développée.